
\subsubsection{4.1 Dataset}

The MBPP benchmark \citep{austin2021mbpp} was used, consisting of 500 Python programming problems in the test split (task IDs 11-510).

\begin{table}[htbp]
\centering
\resizebox{\columnwidth}{!}{%
\begin{tabular}{ll}
\toprule
Property & Value \\
\midrule
Total problems & 500 \\
Split & Test set \\
Problem type & Single-function Python \\
\bottomrule
\end{tabular}
}
\end{table}

\subsubsection{4.2 Procedure}

\textbf{Phase 1 (Error Collection)}: Code was generated for all 500 MBPP problems using CodeLlama-7B-Instruct. Each solution was executed against test cases and categorized as: PASS, RUNTIME\_ERROR, SYNTAX\_ERROR, WRONG\_OUTPUT, or TIMEOUT.

\textbf{Phase 2 (Feedback Generation)}: For each runtime error case, five repair prompts were constructed varying only the error feedback section.

\textbf{Phase 3 (Repair Evaluation)}: Repair was attempted under each condition. Success was recorded as binary: 1 if all tests pass, 0 otherwise.

Total repair attempts: 304 cases x 5 levels = 1,520.

\subsubsection{4.3 Computational Resources}

Experiments were conducted on a single NVIDIA A100 GPU. Code generation for H-E1 required approximately 20 minutes. The full H-M1 repair experiment (1,520 attempts) required approximately 40 minutes.
